% =============================================================================
%  Lab 2 - Master-Slave Flip-Flop and the N-Bit Register
%  Seven sections. Delete the guidance comments as you fill each one in.
% =============================================================================

\chapter{Lab 2 --- Master--Slave Flip-Flop and the N-Bit Register}

\section{Objective}
% Two or three sentences, your own words.

\section{Specification}

\begin{table}[htbp]
\centering
\caption{Transition table (Type T) for \texttt{RegN}.}
\label{tab:lab02-truth}
\begin{tabular}{clll||ll}
\toprule
\# & \texttt{Rst} & \texttt{Ld} & \texttt{Clk} & $Q^{+}$ & note \\
\midrule
0 & 1 & -- & -- & \textbf{0} & asynchronous reset \\
1 & 0 & -- & not $\uparrow$ & \textbf{Q} & hold between edges \\
2 & 0 & 0 & $\uparrow$ & \textbf{Q} & edge, not loading \\
3 & 0 & 1 & $\uparrow$ & \textbf{D} & load \\
\bottomrule
\end{tabular}
\end{table}

\subsection*{The two latch equations}

\begin{align*}
\text{master} \quad & Q_m^{+} = \overline{Clr}\cdot\left( \overline{Clk}\cdot D + Clk\cdot Q_m \right)\\
\text{slave}  \quad & Q^{+}   = \overline{Clr}\cdot\left( Clk\cdot Q_m + \overline{Clk}\cdot Q \right)
\end{align*}

% One sentence: why can Q only change at the rising edge? Argue it from the two
% equations above, not from the phrase "because it is edge-triggered".

\subsection*{The register equation}

\begin{align*}
\text{(1) canonical} \quad & Q_i^{+} = \overline{Rst}\cdot\left( Ld\cdot D_i + \overline{Ld}\cdot Q_i \right)\\
\text{(2) observation} \quad & \text{[\emph{state which Lab 1 signal became which}]}\\
\text{(3) VHDL} \quad & \text{see Listing~\ref{lst:code/lab02/RegN.vhd}}
\end{align*}

\section{Design}
% Two things to answer here:
%   (a) Why does the structural flip-flop exist at all, given that the
%       behavioural one is what you would actually ship?
%   (b) THE SENSITIVITY LIST. Why is RegN's list (Clk, Rst) and not
%       (Clk, Rst, Ld, D)? Take Ld and D one at a time and say why each is
%       absent. This is a graded question.

\section{Implementation}

\vhdlfile{code/lab02/MSFlipFlop.vhd}{\texttt{MSFlipFlop} --- structural, from two transparent latches}

\vhdlfile{code/lab02/RegN.vhd}{\texttt{RegN} --- N-bit register with load enable and asynchronous reset}

\section{Verification}
% Your vectors, then the coverage argument. For this lab the argument has a
% specific shape: which of your vectors would a TRANSPARENT LATCH fail?
% That vector is the one proving the design is edge-triggered.

\begin{table}[htbp]
\centering
\caption{Test vectors. Note \emph{when} each stimulus changes, relative to the clock.}
\label{tab:lab02-vectors}
\begin{tabular}{clllll}
\toprule
\# & when & \texttt{Rst} & \texttt{Ld} & \texttt{D} & expected $Q$ / what it proves \\
\midrule
1 & & & & & \\
2 & & & & & \\
3 & & & & & \\
\bottomrule
\end{tabular}
\end{table}

\section{Results}
% Evidence: E1, E2, E7.

\figplaceholder{the waveform: \texttt{D} toggling three times inside one clock
period while \texttt{Q} moves once --- annotate both}
% \resultfig{figs/lab02/waveform.png}
%   {Between t = ... and t = ..., D changes three times and Q does not move;
%    Q changes only at the rising edge at t = ..., taking the last value D held.}
%   {fig:lab02-wave}

\figplaceholder{RTL Viewer, \emph{structural} version --- two latches and an inverter}

\figplaceholder{RTL Viewer, \emph{behavioural} version --- one flip-flop primitive}

% One sentence on why the two pictures differ despite identical behaviour.

% E4: the Fitter's register count for RegN (top level, N = 32).
%     State the number. It must be 32. Say why that is a correctness check
%     and not merely a cost figure.

\section{Discussion}
% REQUIRED. One bug, how you found it, what it cost.
% If your sensitivity list bit you, that is the bug to write about: say which
% checks failed, and how the pattern of failures told you which signal was
% missing.
